\import{relation/closure.tex}

\section{Transitive closures}

\begin{definition}\label{transitive_supersets}
    $\transitiveSupersets{R} =
        \{S\in\pow{\fld{R}\times\fld{R}}\mid
            \text{$R\subseteq S$ and $S$ is transitive}\}$.
\end{definition}

\begin{definition}\label{transitive_closure}
    $\transitiveClosure{R} = \inters{\transitiveSupersets{R}}$.
\end{definition}

\begin{lemma}\label{field_square_is_transitive}
    $\fld{R}\times\fld{R}$ is transitive.
\end{lemma}
\begin{proof}
    Follows by \auto{transitive,times_tuple_elim,times_tuple_intro}.
\end{proof}

\begin{lemma}\label{transitive_supersets_inhabited}
    Let $R$ be a relation.
    Then $\transitiveSupersets{R}$ is inhabited.
\end{lemma}
\begin{proof}
    $R\subseteq\fld{R}\times\fld{R}$
        by \obv{relation_subseteq_times_fld}.
    $\fld{R}\times\fld{R}\in\pow{\fld{R}\times\fld{R}}$
        by \obv{powerset_intro,subseteq_refl}.
    Thus $\fld{R}\times\fld{R}\in\transitiveSupersets{R}$
        by \obv{transitive_supersets,field_square_is_transitive}.
\end{proof}

\begin{classification}\label{transitive_closure_type}
    Let $R$ be a relation.
    Then $\transitiveClosure{R}$ is a binary relation on $\fld{R}$.
\end{classification}
\begin{proof}
    $\fld{R}\times\fld{R}\in\transitiveSupersets{R}$
        by \obv{transitive_supersets,relation_subseteq_times_fld,powerset_intro,subseteq_refl,field_square_is_transitive}.
    Hence $\transitiveClosure{R}\subseteq\fld{R}\times\fld{R}$
        by \auto{transitive_closure,inters_subseteq_elem}.
    Follows by definition.
\end{proof}

\begin{classification}\label{transitive_closure_binary_relation_on}
    Suppose $R$ is a binary relation on $X$.
    Then $\transitiveClosure{R}$ is a binary relation on $X$.
\end{classification}
\begin{proof}
    $R$ is a relation by notion.
    $\transitiveClosure{R}\subseteq\fld{R}\times\fld{R}$
        by \obv{transitive_closure_type}.
    $\fld{R}\subseteq X$ by \obv{binary_relation_field_subseteq}.
    $\fld{R}\times\fld{R}\subseteq X\times\fld{R}$
        by \obv{times_subseteq_left}.
    $X\times\fld{R}\subseteq X\times X$
        by \obv{times_subseteq_right}.
    Hence $\transitiveClosure{R}\subseteq X\times X$
        by \obv{subseteq_transitive}.
    Follows by definition.
\end{proof}

\begin{proposition}\label{transitive_closure_is_relation}
    Let $R$ be a relation.
    Then $\transitiveClosure{R}$ is a relation.
\end{proposition}
\begin{proof}
    Follows by \auto{relation,transitive_closure_type,times_elem_is_tuple,subseteq}.
\end{proof}

\begin{proposition}\label{transitive_closure_contains}
    Let $R$ be a relation.
    Then $R\subseteq\transitiveClosure{R}$.
\end{proposition}
\begin{proof}
    $\transitiveSupersets{R}$ is inhabited by \obv{transitive_supersets_inhabited}.
    For every $S\in\transitiveSupersets{R}$ we have $R\subseteq S$
        by \obv{transitive_supersets}.
    Follows by \auto{transitive_closure,inters_greatest}.
\end{proof}

\begin{proposition}\label{transitive_closure_is_transitive}
    Let $R$ be a relation.
    Then $\transitiveClosure{R}$ is transitive.
\end{proposition}
\begin{proof}
    $\transitiveSupersets{R}$ is inhabited by \obv{transitive_supersets_inhabited}.
    Follows by \auto{transitive,transitive_closure,inters_iff_forall,transitive_supersets}.
\end{proof}

\begin{lemma}\label{transitive_closure_candidate}
    Let $R$ be a relation.
    Let $S$ be a transitive relation.
    Suppose $R\subseteq S$.
    Then $S\inter(\fld{R}\times\fld{R})\in\transitiveSupersets{R}$.
\end{lemma}
\begin{proof}
    $R\subseteq\fld{R}\times\fld{R}$
        by \obv{relation_subseteq_times_fld}.
    Thus $R\subseteq S\inter(\fld{R}\times\fld{R})$
        by \obv{subseteq_inter_iff}.
    $S\inter(\fld{R}\times\fld{R})\subseteq\fld{R}\times\fld{R}$
        by \obv{inter_lower_right}.
    Hence $S\inter(\fld{R}\times\fld{R})
        \in\pow{\fld{R}\times\fld{R}}$ by \obv{powerset_intro}.
    $S\inter(\fld{R}\times\fld{R})$ is transitive
        by \auto{transitive,inter,times_tuple_elim,times_tuple_intro}.
    Follows by \obv{transitive_supersets}.
\end{proof}

\begin{proposition}\label{transitive_closure_least}
    Let $R$ be a relation.
    Let $S$ be a transitive relation.
    Suppose $R\subseteq S$.
    Then $\transitiveClosure{R}\subseteq S$.
\end{proposition}
\begin{proof}
    $S\inter(\fld{R}\times\fld{R})\in\transitiveSupersets{R}$
        by \obv{transitive_closure_candidate}.
    $\transitiveClosure{R}\subseteq S\inter(\fld{R}\times\fld{R})$
        by \auto{transitive_closure,inters_subseteq_elem}.
    Follows by \auto{subseteq,inter}.
\end{proof}

\begin{proposition}\label{converse_transitive_closure}
    Let $R$ be a relation.
    Then $\converse{\transitiveClosure{R}}
        =\transitiveClosure{\converse{R}}$.
\end{proposition}
\begin{proof}
    $\converse{\transitiveClosure{R}}$ is a transitive relation by
        \obv{transitive_closure_is_relation,
            transitive_closure_is_transitive,converse_is_relation,
            converse_transitive}.
    $\converse{R}\subseteq\converse{\transitiveClosure{R}}$ by
        \obv{transitive_closure_contains,converse_subseteq_intro}.
    Hence $\transitiveClosure{\converse{R}}
        \subseteq\converse{\transitiveClosure{R}}$
        by \obv{transitive_closure_least}.

    $\transitiveClosure{\converse{R}}$ is a transitive relation by
        \obv{transitive_closure_is_relation,
            transitive_closure_is_transitive}.
    $\converse{\transitiveClosure{\converse{R}}}$ is a transitive relation
        by \obv{converse_is_relation,converse_transitive}.
    $R\subseteq\converse{\transitiveClosure{\converse{R}}}$ by
        \auto{subseteq,relation,transitive_closure_contains,
            converse_iff}.
    Hence $\transitiveClosure{R}
        \subseteq\converse{\transitiveClosure{\converse{R}}}$
        by \obv{transitive_closure_least}.
    Therefore $\converse{\transitiveClosure{R}}
        \subseteq\transitiveClosure{\converse{R}}$ by
        \auto{subseteq,relation,converse_iff}.
    Follows by \obv{subseteq_antisymmetric}.
\end{proof}

\begin{proposition}\label{transitive_closure_eq}
    Let $R$ be a transitive relation.
    Then $\transitiveClosure{R}=R$.
\end{proposition}
\begin{proof}
    $R\subseteq\transitiveClosure{R}$ by \obv{transitive_closure_contains}.
    $\transitiveClosure{R}\subseteq R$
        by \obv{transitive_closure_least,subseteq_refl}.
    Follows by \obv{subseteq_antisymmetric}.
\end{proof}

\begin{definition}\label{reflexive_transitive_closure}
    $\reflexiveTransitiveClosure{X}{R}
        = \transitiveClosure{\reflexiveClosure{X}{R}}$.
\end{definition}

\begin{proposition}\label{reflexive_transitive_closure_is_relation}
    Suppose $R$ is a binary relation on $X$.
    Then $\reflexiveTransitiveClosure{X}{R}$ is a relation.
\end{proposition}
\begin{proof}
    $R$ is a relation by \auto{relation,subseteq,times_elem_is_tuple}.
    $\identity{X}$ is a relation by \obv{id_is_relation}.
    Thus $\reflexiveClosure{X}{R}$ is a relation
        by \auto{reflexive_closure,union_relations_is_relation}.
    Follows by \auto{reflexive_transitive_closure,transitive_closure_is_relation}.
\end{proof}

\begin{classification}\label{reflexive_transitive_closure_type}
    Suppose $R$ is a binary relation on $X$.
    Then $\reflexiveTransitiveClosure{X}{R}$ is a binary relation on $X$.
\end{classification}
\begin{proof}
    $\reflexiveClosure{X}{R}$ is a binary relation on $X$
        by \obv{reflexive_closure_binary_relation_on}.
    Hence $\transitiveClosure{\reflexiveClosure{X}{R}}$
        is a binary relation on $X$
        by \obv{transitive_closure_binary_relation_on}.
    Follows by \obv{reflexive_transitive_closure}.
\end{proof}

\begin{proposition}\label{reflexive_transitive_closure_contains}
    Suppose $R$ is a binary relation on $X$.
    Then $R\subseteq\reflexiveTransitiveClosure{X}{R}$.
\end{proposition}
\begin{proof}
    $R$ is a relation by \auto{relation,subseteq,times_elem_is_tuple}.
    $\identity{X}$ is a relation by \obv{id_is_relation}.
    Thus $\reflexiveClosure{X}{R}$ is a relation
        by \auto{reflexive_closure,union_relations_is_relation}.
    $R\subseteq\reflexiveClosure{X}{R}$
        by \auto{reflexive_closure,union_upper_left}.
    $\reflexiveClosure{X}{R}\subseteq\reflexiveTransitiveClosure{X}{R}$
        by \auto{reflexive_transitive_closure,transitive_closure_contains}.
    Follows by \obv{subseteq_transitive}.
\end{proof}

\begin{proposition}\label{reflexive_transitive_closure_is_reflexive}
    Suppose $R$ is a binary relation on $X$.
    Then $\reflexiveTransitiveClosure{X}{R}$ is reflexive on $X$.
\end{proposition}
\begin{proof}
    $R$ is a relation by \auto{relation,subseteq,times_elem_is_tuple}.
    $\identity{X}$ is a relation by \obv{id_is_relation}.
    Thus $\reflexiveClosure{X}{R}$ is a relation
        by \auto{reflexive_closure,union_relations_is_relation}.
    $\reflexiveClosure{X}{R}\subseteq\reflexiveTransitiveClosure{X}{R}$
        by \auto{reflexive_transitive_closure,transitive_closure_contains}.
    Follows by \auto{reflexive_on,reflexive_closure_is_reflexive,subseteq}.
\end{proof}

\begin{proposition}\label{identity_subseteq_reflexive_transitive_closure}
    Suppose $R$ is a binary relation on $X$.
    Then $\identity{X}\subseteq
        \reflexiveTransitiveClosure{X}{R}$.
\end{proposition}
\begin{proof}
    $\reflexiveTransitiveClosure{X}{R}$ is reflexive on $X$
        by \obv{reflexive_transitive_closure_is_reflexive}.
    Show that for every $w\in\identity{X}$ we have
        $w\in\reflexiveTransitiveClosure{X}{R}$.
    \begin{subproof}
        Fix $w\in\identity{X}$.
        Take $x\in X$ such that $w=(x,x)$ by \obv{id_elem_inspect}.
        $x\mathrel{\reflexiveTransitiveClosure{X}{R}}x$
            by \obv{reflexive_on}.
    \end{subproof}
    Follows by \obv{subseteq}.
\end{proof}

\begin{proposition}\label{reflexive_transitive_closure_is_transitive}
    Suppose $R$ is a binary relation on $X$.
    Then $\reflexiveTransitiveClosure{X}{R}$ is transitive.
\end{proposition}
\begin{proof}
    $R$ is a relation by \auto{relation,subseteq,times_elem_is_tuple}.
    $\identity{X}$ is a relation by \obv{id_is_relation}.
    Thus $\reflexiveClosure{X}{R}$ is a relation
        by \auto{reflexive_closure,union_relations_is_relation}.
    Follows by \auto{reflexive_transitive_closure,transitive_closure_is_transitive}.
\end{proof}

\begin{proposition}\label{reflexive_transitive_closure_least}
    Suppose $R$ is a binary relation on $X$.
    Let $S$ be a relation.
    Suppose $S$ is transitive and $S$ is reflexive on $X$.
    Suppose $R\subseteq S$.
    Then $\reflexiveTransitiveClosure{X}{R}\subseteq S$.
\end{proposition}
\begin{proof}
    $R$ is a relation by \obv{relation,subseteq,times_elem_is_tuple}.
    $\identity{X}$ is a relation by \obv{id_is_relation}.
    Thus $\reflexiveClosure{X}{R}$ is a relation
        by \auto{reflexive_closure,union_relations_is_relation}.
    $\identity{X}\subseteq S$
        by \auto{subseteq,relation,id_iff,reflexive_on}.
    Thus $\reflexiveClosure{X}{R}\subseteq S$
        by \auto{reflexive_closure,union_subsets_is_subset}.
    Follows by \auto{reflexive_transitive_closure,transitive_closure_least}.
\end{proof}

\begin{proposition}\label{transitive_closure_subseteq_reflexive_transitive_closure}
    Suppose $R$ is a binary relation on $X$.
    Then $\transitiveClosure{R}\subseteq
        \reflexiveTransitiveClosure{X}{R}$.
\end{proposition}
\begin{proof}
    $\reflexiveTransitiveClosure{X}{R}$ is a transitive relation by
        \obv{reflexive_transitive_closure_is_relation,
            reflexive_transitive_closure_is_transitive}.
    $R\subseteq\reflexiveTransitiveClosure{X}{R}$
        by \obv{reflexive_transitive_closure_contains}.
    Follows by \obv{transitive_closure_least}.
\end{proof}

\begin{theorem}\label{reflexive_transitive_closure_first_step_eq}
    Suppose $R$ is a binary relation on $X$.
    Let $T=\reflexiveTransitiveClosure{X}{R}$.
    Then $T=\identity{X}\union(T\circ R)$.
\end{theorem}
\begin{proof}
    Let $S=\identity{X}\union(T\circ R)$.
    $T$ is a transitive relation and $T$ is reflexive on $X$ by
        \auto{reflexive_transitive_closure_is_relation,
            reflexive_transitive_closure_is_transitive,
            reflexive_transitive_closure_is_reflexive}.
    $R\subseteq T$ by \obv{reflexive_transitive_closure_contains}.
    $S$ is a relation by
        \obv{circ_is_relation,id_is_relation,union_relations_is_relation}.
    $S$ is reflexive on $X$ by
        \auto{reflexive_on,union_iff,id_iff}.
    Show that $R\subseteq S$.
    \begin{subproof}
        Show that for every $w\in R$ we have $w\in S$.
        \begin{subproof}
            Fix $w\in R$.
            Take $x,y$ such that $w=(x,y)$ by \obv{relation}.
            $x\mathrel{R}y$.
            $(x,y)\in X\times X$ by \obv{subseteq}.
            Hence $y\in X$ by \auto{times_tuple_elim}.
            Thus $y\mathrel{T}y$ by \obv{reflexive_on}.
            Hence $x\mathrel{T\circ R}y$ by \auto{circ_iff}.
            Thus $(x,y)\in\identity{X}\union(T\circ R)$
                by \auto{union_iff}.
            Hence $(x,y)\in S$.
        \end{subproof}
        Follows by \obv{subseteq}.
    \end{subproof}
    $S$ is transitive by
        \auto{transitive,union_iff,id_iff,circ_iff,subseteq}.
    Hence $\reflexiveTransitiveClosure{X}{R}\subseteq S$
        by \obv{reflexive_transitive_closure_least}.
    Thus $T\subseteq S$.
    $\identity{X}\subseteq T$ by
        \obv{identity_subseteq_reflexive_transitive_closure}.
    Show that every element of $T\circ R$ is an element of $T$.
    \begin{subproof}
        Fix $w\in T\circ R$.
        Take $x,y,z$ such that $w=(x,z)$ and
            $x\mathrel{R}y\mathrel{T}z$ by \obv{circ}.
        $x\mathrel{T}y$ by \obv{subseteq}.
        Hence $x\mathrel{T}z$ by \obv{transitive}.
    \end{subproof}
    $T\circ R\subseteq T$ by \obv{subseteq}.
    Hence $\identity{X}\union(T\circ R)\subseteq T$
        by \obv{union_subsets_is_subset}.
    Thus $S\subseteq T$.
    $T=S$ by \obv{subseteq_antisymmetric}.
\end{proof}

\begin{theorem}\label{reflexive_transitive_closure_first_step}
    Suppose $R$ is a binary relation on $X$.
    Let $x,y\in X$.
    Then $x\mathrel{\reflexiveTransitiveClosure{X}{R}}y$ iff
        $x=y$ or there exists $z\in X$ such that
            $x\mathrel{R}z$ and
            $z\mathrel{\reflexiveTransitiveClosure{X}{R}}y$.
\end{theorem}
\begin{proof}
    Let $T=\reflexiveTransitiveClosure{X}{R}$.
    $T=\identity{X}\union(T\circ R)$ by
        \obv{reflexive_transitive_closure_first_step_eq}.
    Hence $x\mathrel{T}y$ iff
        $x\mathrel{\identity{X}}y$ or there exists $z$ such that
            $x\mathrel{R}z$ and $z\mathrel{T}y$ by
        \auto{union_iff,circ_iff}.
    $R\subseteq X\times X$ by definition.
    Follows by \auto{id_iff,subseteq,times_tuple_elim}.
\end{proof}

\begin{theorem}\label{converse_reflexive_transitive_closure}
    Suppose $R$ is a binary relation on $X$.
    Then $\converse{\reflexiveTransitiveClosure{X}{R}}
        =\reflexiveTransitiveClosure{X}{\converse{R}}$.
\end{theorem}
\begin{proof}
    $\reflexiveClosure{X}{R}$ is a binary relation on $X$ by notion.
    Hence $\reflexiveClosure{X}{R}$ is a relation by notion.
    Follows by
        \auto{reflexive_transitive_closure,
            converse_transitive_closure,converse_reflexive_closure}.
\end{proof}

\begin{theorem}\label{reflexive_transitive_closure_last_step}
    Suppose $R$ is a binary relation on $X$.
    Let $x,y\in X$.
    Then $x\mathrel{\reflexiveTransitiveClosure{X}{R}}y$ iff
        $x=y$ or there exists $z\in X$ such that
            $x\mathrel{\reflexiveTransitiveClosure{X}{R}}z$ and
            $z\mathrel{R}y$.
\end{theorem}
\begin{proof}
    $\converse{R}$ is a binary relation on $X$
        by \obv{converse_binary_relation_on}.
    Follows by
        \auto{converse_reflexive_transitive_closure,converse_iff,
            reflexive_transitive_closure_first_step}.
\end{proof}
